\paragraph{4. Oracle Construction.}
In this section, we formally present the proposed oracle construction. Let $a>0$ be a parameter.
\paragraph{Construction.}
$\mathcal O_a=\mathcal F+\mathcal A$, where each oracle is randomly selected according to the following process:
\begin{enumerate}
\item Let $t(n)=2^{an/\log n}$ be the upper bound on the time of nonuniform adversaries and $c=7a$.
\item Define functions $p$ and $i_{\max}$ as $p(n)=t(n)^{-6}$ and $i_{\max}(n)=\frac1c\log\log t(n)$. Here, $p$ is a parameter of random restriction, and $i_{\max}$ is the number of applications of random restrictions.
\item For each $n\in\mathbb N$, define a set $V_{n,0}$ of variables taking binary values as follows:
\[V_{n,0}=\{F_{z,x,\ell}:z,x\in\{0,1\}^n,\ell\in[n]\}.\]
\item For each $n\in\mathbb N$, let $S_{n,0}=\{0,1\}^n$.
\item For each $n\in\mathbb N$ and $i\in[i_{\max}(n)-1]$, we inductively (on $i$) select a $p(n)$-random restriction $\rho^*_{n,i}$ to $V_{n,i-1}$ and a random subset $S_{n,i}\subseteq S_{n,i-1}$ of size $p(n)\cdot|S_{n,i-1}|$. Then, we define a restriction $\rho_{n,i}$ to $V_{n,i-1}$ and a subset $V_{n,i}\subseteq V_{n,i-1}$ as follows:
\[\rho_{n,i}(z,x,\ell)=\begin{cases}*&\text{if }z\in S_{n,i},\\\rho^*_{n,i}(z,x,\ell)&\text{otherwise},\end{cases}\]
\[V_{n,i}=\rho_{n,i}^{-1}(*)\left(=\rho_{n,i}^{*-1}(*)\cup\{F_{z,x,\ell}:z\in S_{n,i}\}\right).\]
We also define $\rho_{n,i_{\max}(n)}$ as a full assignment to $V_{n,i_{\max}(n)-1}$ selected uniformly at random. Let $\rho_{n,i}\equiv\rho_{n,i_{\max}(n)}$ for each $i\geq i_{\max}(n)+1$. We use the notation $\rho_{n,\leq i}$ to represent the composite restriction $\rho_{n,1}\cdots\rho_{n,i}$ to $V_{n,0}$ for each $n$ and $i$.
\item Define $\mathcal F=\{F_n\}_{n\in\mathbb N}$, where $F_n:\{0,1\}^n\times\{0,1\}^n\times[n]\to\{0,1\}$, as $F_n(z,x,\ell)=\rho_{n,\leq i_{\max}(n)}(z,x,\ell)$.
\item Define $\mathcal A$ as follows: On input $(M,x,1^{T^{2^c}})$, where $M$ is a nondeterministic oracle machine, $x\in\{0,1\}^*$, and $T\in\mathbb N$, the oracle $\mathcal A(M,x,1^{T^{2^c}})$ returns 0 or 1 according to the following procedure:
\begin{enumerate}
\item Let $i_T:=\frac1c\log\log T$.
\item Construct a $T$-DNF $\phi$ on variables in $V_{n,0}$ representing the execution of $M^{\mathcal F+\mathcal A}(x)$ in $T$ steps, where the top-most OR corresponds to the nondeterminism on a possible choice of $\mathcal F$ (say, $F'$) and an accepting path of $M^{F'+\mathcal A}(x)$, and each term performs verification whether $M$'s at most $T$ queries (say, $(z_1,x_1,\ell_1),\ldots,(z_q,x_q,\ell_q)$ for some $q\leq T$) are consistent with the actual choices of $\mathcal F$, i.e., for each $i\in[q]$, the term contains $F_{z_i,x_i,\ell_i}$ if $F'(z_i,x_i,\ell_i)=1$; otherwise, $\neg F_{z_i,x_i,\ell_i}$ as a literal.
\item If $\phi|_{\rho_{1,\leq i_T},\ldots,\rho_{T,\leq i_T}}\equiv b$ for some $b\in\{0,1\}$, then return $b$, otherwise, return 0.
\end{enumerate}
\end{enumerate}
We can verify that $\mathcal A$ is well-defined (i.e., not circular on recursive calls for $\mathcal A$) as follows.
\paragraph{Proposition 14.}
For each input, the value of $\mathcal A(M,x,1^{T^{2^c}})$ is determined only by $\rho_{n,j}$ (equivalently, $\rho^*_{n,j}$ and $S_{n,j}$) for $n\leq T$ and $j\leq i_T$ (remember that $i_T=\frac1c\log\log T$).
\begin{proof}
We show the proposition by induction on $T$. Remember that, on input $(M,x,1^{T^{2^c}})$, the oracle $\mathcal A$ first makes a $T$-DNF $\phi$ based on $M$ independently of the values of $\mathcal F$.
Suppose that $M$ makes some query $(M',x',1^{T'^{2^c}})$ to $\mathcal A$ for constructing $\phi$. Since the length of such a query is at most $T$, we have $T'^{2^c}\leq T$ and
\[i_{T'}=\frac1c\log\log T'\leq\frac1c\log\log T^{1/2^c}=\frac1c\log\log T-1=i_T-1.\]
By the induction hypothesis, the answer of $\mathcal A(M',x',1^{T'^{2^c}})$ is determined by only $\rho_{n,j}$ for $n\leq T'$ and $j\leq i_T-1$, and so is $\phi$. Then, $\mathcal A$ determines the answer by restricting $\phi$ by $\rho_{n,j}$ for $n\leq T$ and $j\leq i_T$. Therefore, $\mathcal A(M,x,1^{T^{2^c}})$ is determined only by $\rho_{n,j}$ for $n\leq T$ and $j\leq i_T$.
\end{proof}
When the parameter $a$ is clear from the context, we omit the subscript $a$ from $\mathcal O_a$.
